\begin{textsample}{POS Dim 2 – ma – Score 41.00 – ma/ma\_em\_53.txt}  \label{ex:f2_pos_051}
Princípios norteadores do \textbf{itinerário} de formação \textbf{técnica} e \textbf{profissional} A LDB nº 9.394/96 \textbf{preconiza} a educação \textbf{profissional} e tecnológica ( EPT ) como modalidade educacional cujo princípio basilar consiste na integração aos diversos níveis e modalidades de educação, às dimensões do trabalho, da ciência e da tecnologia, visando contribuir para que o cidadão possa se inserir e atuar plenamente no mundo do trabalho e na vida em sociedade.
Os níveis e modalidades da educação \textbf{profissional} e \textbf{técnica}, em articulação com a Educação de Jovens e Adultos ( EJA ), Educação Especial, Educação Escolar Quilombola, Educação Escolar Indígena, Educação do Campo, Educação a Distância ( EaD ) e Educação de Pessoas em \textbf{Regime} de Acolhimento e Privação de Liberdade como possibilidades de aprofundamentos de conhecimentos que devem estar ao alcance do estudante, no âmbito da \textbf{rede} de ensino do estado do Maranhão, poderão ser realizadas em instituição próprias de ensino ou em parceria com instituições credenciadas pelo Conselho Estadual de Educação e \textbf{homologadas} pelo \textbf{Secretário} Estadual de Educação, além de observadas as \textbf{regulamentações} específicas a cada modalidade de ensino, entre outras, a saber : - a Resolução nº 1, de 28 de maio de 2021, que \textbf{institui} \textbf{Diretrizes} Operacionais para a Educação de Jovens e Adultos nos aspectos relativos ao seu alinhamento à Política Nacional de Alfabetização ( PNA ) e à Base Nacional Comum Curricular ( BNCC ), e Educação de Jovens e Adultos a Distância, destacando, no art. 3º, inciso III, que o Ensino médio deve objetivar uma formação geral básica e \textbf{profissional} mais consolidada, tanto na \textbf{oferta} integrada a uma \textbf{qualificação} \textbf{profissional} quanto na \textbf{oferta} de \textbf{curso} \textbf{técnico} de nível médio, com \textbf{carga} \textbf{horária} total mínima de 1.200 horas. - o Parecer CNE/CP nº 22/2020, aprovado em 8 de dezembro de 2020, que define \textbf{Diretrizes} Curriculares da Pedagogia da \textbf{Alternância} na Educação Básica e na Educação Superior.
Este documento apresenta a Pedagogia da \textbf{Alternância} como potencial para organização do processo de formação em atendimento às demandas educacionais de parcelas significativas da população brasileira, cujos princípios \textbf{amparam} a integração de conhecimentos práticos, científicos, diversidade epistemológica, identidades, saberes, territórios educativos e territorialidades dos sujeitos no âmbito da escola, da universidade e de outras instituições educacionais.
Como projeto \textbf{formativo}, a Pedagogia da \textbf{Alternância} responde aos interesses das comunidades urbanas, sobretudo, das comunidades indígenas, quilombolas e povos tradicionais. - a Resolução nº 2, de 11 de setembro de 2021, ao \textbf{instituir} \textbf{Diretrizes} Nacionais para a Educação Especial na Educação Básica, delibera que as escolas que \textbf{ofertam} educação \textbf{profissional} ( públicas e privadas ) devem \textbf{atender} alunos com necessidades educacionais especiais, mediante : promoção das condições de acessibilidade ; \textbf{capacitação} de recursos humanos ; e flexibilização e adaptação do \textbf{currículo} e o encaminhamento para o trabalho ( no Art. 17 ).
Podendo também avaliar e certificar competências laborais de pessoas com necessidades especiais não matriculadas em seus \textbf{cursos}, de modo a conduzi -las para o mundo do trabalho.
Nesse sentido, a \textbf{oferta} de Educação Profissional e Tecnológica ( EPT ), no que concerne aos objetivos, especificidades e duração dos níveis e modalidades, observará as orientações previstas na Lei nº 13.415/2017, nas DCNEM/2018, nas normas do Sistema de Ensino Estadual, do Conselho Estadual do Maranhão, especialmente a Resolução 277/2021-CEE/MA, que estabelece normas complementares para implementação do ensino médio nos termos da Lei nº 13.415/2017, no âmbito do Sistema Estadual de Ensino do Estado do Maranhão, assim como demais normas \textbf{vigentes} que direcionam a \textbf{oferta} da educação \textbf{profissional} e \textbf{técnica} no país, como a Resolução nº 2, de 15 de dezembro de 2020, que aprova a quarta edição do \textbf{Catálogo} Nacional de \textbf{Cursos} \textbf{Técnicos}.
A \textbf{oferta} de Educação Profissional e Tecnológica prevista neste documento concebe uma formação educacional que caminha para o atendimento \textbf{precípuo} dos interesses dos estudantes e das demandas das comunidades local e regional, delineadas em uma organização curricular consolidada nos planos de \textbf{cursos}, propostas pedagógicas institucionais e, consequentemente, no cotidiano escolar.
Nessa perspectiva, a \textbf{oferta} da Educação Profissional e Tecnológica na \textbf{rede} \textbf{estadual} do Maranhão será orientada pelos princípios estabelecidos nas \textbf{legislações} acima mencionadas, com ênfase nas Resoluções CNE/CP nº 1, de 5 de janeiro de 2021, e nº 3/2018 ( DCNEM ), quais sejam : - formação integral do estudante a partir da construção curricular, articulando formação geral básica à formação \textbf{profissional} específica, manifestada por valores, aspectos físicos, cognitivos e socioemocionais, em busca da superação da \textbf{fragmentação} de conhecimentos e segmentação da estruturação curricular ; - construção de percursos \textbf{formativos} que concebam o projeto de vida dos estudantes como estratégia fundamental para o desenvolvimento das dimensões pessoal, cidadã e \textbf{profissional} ; - prática da pesquisa como princípio pedagógico para inovação, criação, resolução de problemas mais complexos e construção de novos conhecimentos ; - trabalho como princípio educativo e base para a organização curricular, com fins de construção das competências \textbf{profissionais}, em seus objetivos, conteúdos e estratégias de ensino e aprendizagem, visando a sua integração com a ciência, a cultura e a tecnologia ; - efetivo desenvolvimento da pessoa, seu preparo para o exercício da cidadania e sua \textbf{qualificação} para o trabalho a partir de uma formação focada no respeito aos direitos humanos como direito universal para o \textbf{fortalecimento} de valores estéticos, políticos e éticos da educação nacional ; - reconhecimento e respeito à diversidade, às identidades dos sujeitos brasileiros, ao pluralismo de ideias, culturas e concepções pedagógicas, bem como das variadas formas de produção e de trabalho, com vistas à inclusão educacional e social ; - conexão entre as múltiplas dimensões do eixo tecnológico que norteia a \textbf{oferta} de um \textbf{curso} ao conhecimento científico e tecnológico vinculados à \textbf{oferta}, tendo como fundamento a sustentabilidade ambiental ; - flexibilização na construção de percursos \textbf{formativos} como forma de possibilitar múltiplas \textbf{trajetórias} aos estudantes, a partir da articulação entre saberes e os contextos histórico, econômico, social, científico, ambiental, cultural e do mundo do trabalho ; - reconhecimento da historicidade dos conhecimentos e das experiências de vida dos estudantes na estruturação curricular, fundada na indissociabilidade entre teoria e prática, entre educação e prática social ; - \textbf{oferta} de \textbf{cursos} articulada com os arranjos produtivos e demandas locais e regionais, visando ao desenvolvimento socioeconômico e ambiental dos territórios ; - autonomia das instituições de ensino na elaboração e reelaboração de suas propostas pedagógicas de forma democrática e participativa, em atendimento às \textbf{legislações} e às normas educacionais nacionais e \textbf{estaduais} ; - construção do \textbf{perfil} \textbf{profissional} de conclusão de \textbf{curso} \textbf{consubstanciado} pela interrelação de conhecimentos gerais e específicos, competências e habilidades do mundo do trabalho, do desenvolvimento tecnológico e das demandas sociais, econômicas e ambientais ; - \textbf{fortalecimento} do \textbf{regime} de colaboração entre entes federados, visando ao estabelecimento de parcerias para ampliação de possibilidades de \textbf{oferta} aos estudantes, com \textbf{equidade} e \textbf{qualidade} educacional.
As atuais relações de trabalho são construídas sobre os alicerces do sujeito \textbf{trabalhador} global, \textbf{flexível}, criativo, que detém o todo no processo produtivo.
Para tanto, esse \textbf{trabalhador} necessita de uma formação que lhe possibilite a construção de conhecimentos técnico-científicos articulados ao desenvolvimento intelectual e socioemocional, contemplando um saber que não se encerra no processo manual, aplicável, mecânico, típico de formação para um trabalho simples, mas em uma formação integral, que incorpora fundamentos do trabalho, das práticas pessoais, sociais, culturais, científicas e tecnológicas.
Aspectos de extrema relevância para a construção de um \textbf{perfil} \textbf{profissional} capaz de desenvolver competências e habilidades específicas de determinada área \textbf{profissional}, de forma que os estudantes possam atuar com autonomia e criticidade, sendo protagonistas da sua história de vida e \textbf{profissional}.
Nessa direção, as configurações de \textbf{ofertas} do \textbf{itinerário} de formação \textbf{profissional} e \textbf{técnica} compreendem terminologias e conceitos próprios estabelecidos pela Resolução nº 3/2018 ( DCNEM ), que estão especificados no quadro a seguir.
Quadro 4 – Itinerário \textbf{profissional} e \textbf{técnico} : termos e conceitos próprios | Termo | Conceito | |---|---| | Ambiente simulado | Ambientes pedagógicos para atividades práticas da aprendizagem \textbf{profissional}. | | Formações experimentais | Formações autorizadas pelos respectivos sistemas de ensino que ainda não constam do \textbf{Catálogo} Nacional de \textbf{Cursos} \textbf{Técnicos}. | | Aprendizagem \textbf{profissional} | Formação compatível com o desenvolvimento físico, moral, psicológico e social do jovem na faixa etária de 14 a 24 anos de idade. | | \textbf{Qualificação} \textbf{profissional} | Processo de formação com o desenvolvimento de competências de um \textbf{perfil} \textbf{profissional} demandado pelo mercado de trabalho. | | \textbf{Habilitação} \textbf{profissional} \textbf{técnica} de nível médio | Formação \textbf{profissional} reconhecida por meio de diploma de conclusão de \textbf{curso} \textbf{técnico} com validade nacional. | | Programa de aprendizagem | Articulação de \textbf{cursos} que possibilitam aproveitamentos curriculares por \textbf{itinerários} \textbf{formativos}. | | \textbf{Certificação} intermediária | \textbf{Certificação} de \textbf{qualificação} para o trabalho associado à estruturação curricular por etapas com terminalidade. | | \textbf{Certificação} \textbf{profissional} | Processo de avaliação, reconhecimento e \textbf{certificação} de conhecimentos advindos da educação \textbf{profissional}, inclusive desenvolvida no trabalho. | Fonte : Organizado pelos autores a partir de BRASIL, 2018

% matched lemmas: alternância, amparar, atender, capacitação, carga, catálogo, certificação, consubstanciar, currículo, curso, diretriz, equidade, estadual, flexível, formativo, fortalecimento, fragmentação, habilitação, homologar, horário, instituir, itinerário, legislação, oferta, ofertar, perfil, preconizar, precípuo, profissional, qualidade, qualificação, rede, regime, regulamentação, secretário, trabalhador, trajetória, técnico, vigente
\end{textsample}
